\section{Full permanent multiplicities from coupled shared cores}

\subsection*{Finding F-023: the even-order permanent jet map}
Let \(r\geq4\), \(\ell\geq1\), and define
\[
 R=r-1,\qquad b=2\ell+2,\qquad h=\lceil\sqrt r\rceil,
 \qquad L=2\binom{r+\ell-1}{\ell}-r,\qquad N=hL.
\]
If \(n\geq m\geq N\), then every partition
\(\lambda\vdash nd\) with \(\ell(\lambda)\leq r\)
and \(\lambda_2\leq b\) satisfies
\begin{equation}\label{eq:F023full}
           q^m_\lambda(d[n])=K_n(\lambda)=a_\lambda(d[n]).
\end{equation}
Here the permanent is internally padded by
\(X_{11}^{n-m}\). The equality includes every multiplicity
copy and every outer degree \(d\). Its geometric input is an
integrally split differential of the permanent coefficient-jet
map, with the distinguished matrix entry held equal to the
first variable. All representation decompositions below are
in characteristic zero.

\paragraph{A positive graph and its root cycle covers.}
Use exactly the \(h\) coupled shared-core components of
M-016: their local graphs, roots \(o_a\), monomial trees,
cyclic permutation \(\sigma\), and root bridge matrix
\(C=(C_{ac})\), with \(C_{hh}=0\). All nonzero labels
are positive variables. Let \(B(z)\) be the adjacency
matrix and take \(A_+=xI+B(z)\), where
\(z=(z_1,\ldots,z_R)\). Put
\[
 s=\sum_{i=1}^R z_i z_{\sigma(i)},\qquad
                       P=I_h+C+sI_h.
                                                               \tag{F-023.1}
\]
We claim
\[
                      \operatorname{per}A_+(1,z)
                                  =\operatorname{per}P.
                                                               \tag{F-023.2}
\]
Indeed every nontrivial cycle either is a local two-edge
star cycle through one root, or is a cycle of bridges among
roots. Higher inward and outward tree vertices belong to
no cycle. In a cycle cover, a root not used by a bridge
cycle can use its diagonal weight \(1+C_{aa}\), or choose
one local star cycle, whose total weight is \(s\).
All unused nonroot vertices use their diagonal loops.
These choices give precisely the permutation expansion of
\(\operatorname{per}P\). This is a positive cycle-cover
identity over \(\mathbb Z\), without a permanent version
of Schur-complement elimination.

Define the oriented root minors
\[
 Q_{ac}=\operatorname{per}P_{\widehat c,\widehat a},
                                                               \tag{F-023.3}
\]
where the subscript means deleting row \(c\) and column
\(a\). For \(a\ne c\), this is the minor whose open
path runs from root \(a\) to root \(c\). The first terms
of their permutation expansions are
\[
 \begin{split}
 Q_{ac}&=C_{ac}+O(z^2)\qquad(a\ne c),\\
 Q_{aa}&=1+\operatorname{tr}C-C_{aa}+O(z^2),\\
 Q_{hh}-Q_{aa}&=C_{aa}+O(z^2),\qquad
 Q_{hh}=1+O(z).
 \end{split}                                             \tag{F-023.4}
\]
For the second line the linear terms are the choices of one
nonconstant diagonal entry in the remaining roots; an off-diagonal
permutation cycle uses at least two bridges. For the first line,
the sole degree-one open path is the bridge \(a\to c\).
Thus every displayed leading coefficient is exactly one.

\paragraph{The selected full-matrix minors.}
For every monomial \(z^\gamma\) of degree at most
\(2\ell\), F-022 provides, in each component \(a\),
a local simple path of that weight which uses its root.
The degree-zero case means varying the root diagonal.
For positive degree the endpoints are distinct; for a path
with both endpoints nonroot, the shared core vertices at
its inward and outward ends are distinct. This removes
any path bypassing the root.

The corresponding full-matrix permanent cofactor is exactly
\[
                                  z^\gamma Q_{aa}.
                                                               \tag{F-023.5}
\]
To see this directly, a cofactor permutation has one open
path between the chosen endpoints and disjoint cycles.
The local path occupies root \(a\), so that root cannot
be revisited or belong to another cycle. A path leaving and
returning to that root would revisit it and is forbidden.
The remaining choices at all other roots are exactly the
cycle covers of \(P_{\widehat a,\widehat a}\).
All remaining nonroot vertices of component \(a\) are
forced to use diagonal loops. This proves (F-023.5).

For distinct components \(a,c\), take an inward monomial
vertex of weight \(z^\alpha\) in component \(a\) and
an outward one of weight \(z^\beta\) in component \(c\),
with \(|\alpha|=|\beta|=\ell\). The full-matrix
cofactor is
\[
                                  z^{\alpha+\beta}Q_{ac}.
                                                               \tag{F-023.6}
\]
Here the open path starts with the unique local inward path,
uses an open root path \(a\to c\), and ends with the
unique local outward path. It cannot bypass a root because
the endpoints belong to distinct components. On every root
used by that path the local star cycles are unavailable;
on the other roots the same choices as in (F-023.2) remain.
The root open path and its disjoint cycles are exactly the
minor expansion (F-023.3). With the adjacency convention
\(B_{uv}\) for \(u\to v\), the varied matrix entry
is the reverse entry \((v,u)\). Its cofactor deletes row
\(v\) and column \(u\); this explains the deletion order
in (F-023.3) and fixes the orientation of (F-023.6).

\begin{proof}[Integral jet dominance with the padding entry fixed]
For a target monomial of degree \(b=2\ell+2\), remove
one occurrence as a bridge label and another as the linear
entry variation \(H\). The remaining monomial has degree
\(2\ell\). If the bridge occupies an off-diagonal slot
\((a,c)\), split that remainder into two degree-\(\ell\)
monomials. Equations (F-023.6) and (F-023.4), multiplied
by \(H\), give exactly the target monomial modulo terms
of degree greater than \(b\).

If the bridge occupies a diagonal slot \((a,a)\), then
\(a\ne h\). Use the same remainder \(z^\gamma\)
in components \(h,a\), and subtract their two elementary
directions in that order. By (F-023.5), their difference
is \(H z^\gamma(Q_{hh}-Q_{aa})\), whose truncation
is exactly the target monomial. Thus the entire top degree
is in the image of the differential over \(\mathbb Z\).

For every lower degree at most \(2\ell\), use
(F-023.5) in component \(h\), varied by \(H=x\).
The resulting cofactor has the required monomial as its
lowest term, with coefficient one. Degree \(2\ell+1\)
comes from a degree-\(2\ell\) cofactor times a variable.
Descending by degree and subtracting the higher monomials
already obtained gives every target monomial exactly.
Only integral sums and differences are used. This constructs
an integral linear right inverse to the differential.

Hold one nonroot core diagonal, for instance the diagonal
at \(v_1\) in the first component, fixed as \(x\).
None of these selected directions varies it. The constant
cofactor uses a root diagonal, and all positive-degree
local and cross-component paths have distinct endpoints.
Their elementary variations are off diagonal. The diagonal
comparison in the top-degree argument compares two such
off-diagonal cofactor directions, rather than varying a
nonroot diagonal. Simultaneously permute rows and columns
to place the fixed entry at \((1,1)\).

The order of the jet is allowed: F-022 gives
\(L\geq2\ell+1\), and \(h\geq2\), so \(N\geq b\).
For \(m>N\), append an \(xI_{m-N}\) block. Restrict
the source to the integral affine space specified by the
right-inverse directions and translate the target by its
base coefficient vector. The resulting polynomial map is
\(F(t)=t+O(t^2)\) over \(\mathbb Z\). The lowest
nonzero homogeneous term of any polynomial survives this
substitution, so pullback is injective after every commutative
base change. In particular the jet projection is dominant
over \(\mathbb C\).

Substitution into the internally padded permanent is exactly
\(x^{n-m}\operatorname{per}A_+\), since its distinguished
entry is fixed as \(x\). After \(x=1\), this has the
same detected coefficients. Every such substituted form lies
in the permanent orbit closure. A highest-weight ideal vector
with \(\lambda_2\leq b\) is a polynomial in these jet
coordinates by F-016. Injective pullback forces it to vanish
as a polynomial. Hence no ambient highest-weight copy is lost
on the permanent closure, proving \(q^m_\lambda=a_\lambda\).
M-016 gives \(K_n=a_\lambda\) because \(n\geq m\geq N\).
This proves \eqref{eq:F023full}. If the ambient multiplicity
is zero, the same conclusion holds automatically.
\end{proof}

\paragraph{A smaller explicit positive family beyond every exact cover.}
At \(r=1024\), \(b=4\), one has \(h=32\), \(L=r\),
and \(N=32768\). Take
\[
 m=32768,\quad n=32769,\quad d=2048=2r,\qquad
                  \lambda=(nd-4(r-1),4^{r-1}).
                                                               \tag{F-023.7}
\]
The full nested coefficient support is the degree-four
simplex in \(R=1023\) variables. The checked double-counting
bound \eqref{eq:F021cover} in F-021 gives
\[
          \kappa(\Gamma)\geq1+\left\lceil
                         \frac{1023\cdot1022}{12}\right\rceil
                      =87127>n>m.
\]
Thus every exact two-tree support cover of F-018 fails its
size condition, while \eqref{eq:F023full} applies.
The two independent highest-weight polynomials constructed
in F-021---the square of the leading Hessian determinant,
and \(I(F)F(e_1)^r\) from F-011 at \(k=4\)---have
the displayed weight for these values of \(n,d\).
The first is not divisible by \(F(e_1)\), whereas the
second is; their proof of independence is unchanged.
Consequently this example satisfies
\[
                              2\leq q^m_\lambda=K_n=a_\lambda.
\]
The comparison is with every support cover in that explicit
architecture. It is not a lower bound for arbitrary
representations of the forms and does not provide an
actual obstruction.

\paragraph{Grothendieck dictionary.}
The graph, cycle-cover contractions, root minors and affine
coefficient maps are defined over \(\mathbb Z\). The jet
module is finite locally free, and the integral differential
splitting yields an injection of coordinate algebras universal
under base change. Torus weights and ideal comodules are
those of SGA~3, I, Propositions~4.7.2--4.7.3, pp.~34--36.
The passage to all highest-weight copies is F-016's explicit
Schur-functor support calculation, using SGA~6, VI,
Theorem~3.3, pp.~392--393, and Macdonald, I(4.3), p.~63,
Appendix~A(5.3)--(5.6), pp.~155--156.
An equality in representation \(K_0\) alone supplies
neither the coordinate injection nor the fixed internal
padding entry; both are established by the actual integral
morphisms here. The result determines full permanent and
determinant multiplicities in the stated range, without
claiming a general saturation principle or a comparison
outside that range.
